\chapter*{Introduzione}

Questa tesi di Laurea triennale in Matematica si pone come obiettivo quello di dimostrare un ben noto risultato della geometria differenziale, ovvero il seguente teorema. 
\paragraph{Teorema} Se $ X $ e $ Y $ sono due campi vettoriali di classe $ C^{2} $ su una varietà differenziabile $ M $, i loro flussi commutano se e solo se $ [X,Y] \equiv 0 $.

L'approccio scelto utilizza strumenti dell'analisi non-standard. In particolare, nei capitoli 2 e 3 si fornisce una rigorosa trattazione della costruzione dell'estensione non-standard di $ \R $ ed una generalizzazione all'universo $ V(\R) $. Nel capitolo 4, infine, si utilizzano le basi teoriche poste in precedenza per definire l'estensione non standard di una varietà differenziabile $ M $, e per definire la naturale traduzione dei campi vettoriali in ambito infinitesimale. 

Con queste premesse, il flusso di un campo vettoriale, inizialmente definito come oggetto ``integrale'', diventa un cammino composto da passi infinitesimali. Questo consente una definizione formale di flusso che concordi con la intuizione naturale del concetto.